\section{Type IIB and its Puzzles}
We set the stage for the type IIB supergravity by writing down its spectrum, action and puzzles.
Early on, it was understood that to have a unitary, interacting theory of gravity, the particles can have at most spin 2 \cite{Weinberg:1964ew,Weinberg:1965nx}, which limits the amount of supercharges to be at most 32 \cite{Nahm:1977tg}. 
The most notable theories with 32 supercharges are the 11d, type IIA, and type IIB supergravities. 
These theories are believed to be the low-energy limits of the corresponding membrane and string theories \cite{Green:2012oqa}.

{\color{orange}
Einstein's theory of gravity is nonlinear with an infinite-dimensional diffeomorphism symmetry group.
However, one can expand it to quadratic order around a flat spacetime with $g_{mn} = \eta_{mn} + \kappa h_{mn}$, where $\kappa$ is the gravitational coupling.
Then the Einstein-Hilbert action becomes that of a free massless spin-2 field $h_{mn}$, with a rigid Poincare symmetry SO(d-1,1), as well as possibly a rigid R-symmetry acting on the supercharges.
On-shell single-particle (free) states furnish representations of the little group SO(d-2), From this one can determine the theory's low-energy spectrum.

The upper-bound of 32 supercharges implies that a supergravity theory can have at most 128 bosonic and 128 fermionic on-shell degrees of freedom.
In 11d, they transform under irreducible representations of SO(9), and consists of a gravitino, a graviton, and a 3-form potential.
The 11d supergravity multiplet can be given with SO(9) Dynkin labels as\footnote{Our convention for Dynkin Labels is given in Appendix \ref{sec:Dynkin Label Conventions}.}
\begin{equation}
G_{11} = [2000]_{SO(9)} + [0010]_{SO(9)} + [1001]_{SO(9)}.
\label{11d SUGRA multiplet}
\end{equation}
This can be reduced to 10d by decomposing SO(9) representations into those of SO(8)\footnote{
Specifically, the decomposition is given by branching rules of the form
\begin{equation}
\begin{aligned}
[2000]_{SO(9)} &= [2000]_{SO(8)} + [1000]_{SO(8)}+[0000]_{SO(8)}\\
[0010]_{SO(9)} &= [0011]_{SO(8)} + [0100]_{SO(8)}\\
[1001]_{SO(9)} &= [1001]_{SO(8)} + [1010]_{SO(8)}+[0010]_{SO(8)}+[0001]_{SO(8)}.
\end{aligned}
\end{equation}
}.
One finds the type IIA multiplet (with SO(8) Dynkin labels)
\begin{equation}
\begin{aligned}
G_{IIA} 
&= \underbrace{[2000]+[0100]+[0000]}_{NSNS} + \underbrace{[0011] + [1000]}_{RR} + [1001] + [0010] + [1010] + [0001]\\
&=([1000] + [0001]) \times ([1000] + [0010]).
\end{aligned}
\label{IIA multiplet}
\end{equation}
}
Remarkably, the type IIA multiplet \eqref{IIA multiplet} factorizes, as the product of two vector multiplets of different chirality in 10d.
In this product, the bosonic fields are classified into NSNS and RR sectors: the NSNS sector is that obtained by tensor product between two vectors, and the RR sector is that obtained by two spinors \cite{Polchinski:1998rr}.
We note that the NSNS/RR distinction is not apparent when the IIA multiplet is viewed as dimensionally reduced from 11d supergravity, but only becomes distinguished when viewed as 10d tensor products.
This is related to how membranes live naturally in 11d \cite{Bergshoeff:1987cm}, while strings live in 10d \cite{Green:2012oqa}, and the NSNS/RR classification has a stringy origin \cite{Polchinski:1995mt}.

The factorization of the type IIA multiplet \eqref{IIA multiplet} leads one to construct the other supergravity multiplet with 32 supercharges, by instead taking the tensor product of two vector multiplets of the same chirality. This is the (chiral) type IIB supergravity multiplet
\begin{equation}
\begin{aligned}
G_{IIB}
&=\left([1000] + [0001]\right)^2\\
&=\underbrace{[2000]+[0100] + [0000]}_{NSNS} 
+ \underbrace{[0002] + [0100] + [0000]}_{RR} + 2\cdot[1001] + 2\cdot[0010].
\end{aligned}
\end{equation}
Unlike the type IIA supergravity theory, the type IIB theory in 10 dimensions is not known to follow from dimensional reduction of another.
This is the first puzzle of the type IIB theory: does it have a higher dimensional origin?

{\color{orange}
Interestingly, the type IIB multiplet can only contain the components of the 4-form potential $C_4$ whose field strength is (anti)self-dual. 
This is demanded by supersymmetry as had we included the full 4-form potential, the bosonic degrees of freedom would exceed 128.}
This leads to the second major puzzle of the type IIB theory: 
what is the dynamical mechanism behind self-duality of the 5-form field strength?
This can not be imposed at the level of the action, e.g. via Lagrange multipliers\footnote{
A simple way to see this is to start with a generic 4-form $C_4$ with field strength $F_5$, which has 70 on-shell degrees of freedom in 10d.
Adding a Lagrange multiplier $\Lambda_4$ imposing $F_5 = * F_5$.
The on-shell degrees of freedom now contains two 4-forms, the self-duality constraint only removes half, leaving again 70 degrees of freedom.
}. 
In practice, one imposes the self-duality as an additional equation or follow the PST formalism \cite{Pasti:1996vs, DallAgata:1997gnw, DallAgata:1998ahf}, which allows one to derive the self-duality condition from an action, by introducing extra scalar fields along with gauge invariance.

{\color{orange}
We have been discussing the type IIB supergravity multiplet, which captures the massless spectrum of the quadratic (free) theory with SO(9, 1) symmetry, and an SO(2) $\cong$ U(1) rigid symmetry acting on the supercharges.
The nonlinear (interacting) theory is obtained by restoring an infinite number of higher order terms in the gravitational coupling $\kappa$, which promote the rigid SO(9,1) to diffeomorphism
\footnote{
For example, promoting $\frac{1}{2}\int d^{10}x\left(\frac{1}{2}\partial_p h_{mn}\partial^p h^{mn}-\frac{1}{4}\partial_m h \partial^m h\right)$ in de Donder gauge to $\frac{1}{2\kappa^2}\int d^{10}x\sqrt{-g}R$}
and rigid U(1) to local U(1). Doing so \cite{Green:1982tk,Schwarz:1983wa,Schwarz:1983qr,Howe:1983sra} one finds that the theory has a hidden global symmetry group SL(2, R), which when combined with the now localised U(1) symmetry, gives $\frac{SL(2, R)}{U(1)}$ as the scalar manifold.



The action of the type IIB supergravity \cite{Schwarz:1983qr,Howe:1983sra,Bergshoeff:1995as} with the string frame metric $G_{mn}$ can be written as 
\cite{Becker:2006dvp,ValeixoBento:2025emh}
\footnote{
We will not follow the PST approach, and instead impose self-duality as an additional field equation.
}
\begin{equation}
\begin{aligned}
S_{IIB}
&=S_{NSNS} + S_{RR} + S_{CS}\\
&=\frac{1}{2\kappa^2}
\int d^{10}x \sqrt{-G}
\Bigg[
e^{-2\Phi}\left(R + 4(\partial\Phi)^2-\frac{1}{2}|H_3|^2\right)
-\left(\frac{1}{2}(\partial C)^2+\frac{1}{2}|\tilde F_3|^2+\frac{1}{4}|\tilde{F}_5|^2\right)
\Bigg]\\
&\quad -\frac{1}{4\kappa^2}\int C_4\wedge H_3\wedge F_3
\end{aligned}
\label{string frame IIB action}
\end{equation}
with $2\kappa^2=(2\pi)^7l_s^8=(2\pi)^7\alpha^{\prime 4}$, and
\begin{equation}
F_p = dC_{p-1}, \quad 
H_3 = dB_2, \quad
\tilde F_3 = F_3 - C H_3,\quad
\tilde F_5 = * \tilde F_5 = F_5 - \frac12 C_2\wedge H_3 + \frac12 B_2\wedge F_3.
\end{equation}
We note that \eqref{string frame IIB action} is invariant, up to a boundary term, under the SL(2, R) transformations
\begin{equation}
G_{mn}\to |c\tau+d|G_{mn},\quad
\tau\equiv C + i e^{-\Phi} \to \frac{a\tau+b}{c\tau+d},\quad
\begin{pmatrix}
F_3\\H_3
\end{pmatrix}
\to
\begin{pmatrix}
d&c\\b&a
\end{pmatrix}
\begin{pmatrix}
H_3\\F_3
\end{pmatrix},\quad
F_5 \to F_5.
\label{string frame sl2r transformations}
\end{equation}

The presence of SL(2, R) acting on the metric \eqref{string frame sl2r transformations} is a consequence of the fact that the string frame Einstein-Hilbert term couples to the dilaton.
The standard procedure is to then define an Einstein frame metric \cite{Becker:2006dvp,Weigand:2018rez} where the Einstein-Hilbert action is canonical.
This fixes the Einstein frame metric up to a conventional factor.
When it comes to fixing such factor, two choices are common, either by
\begin{enumerate}
    \item demanding that the Einstein frame metric is SL(2, R)-invariant, or
    \item demanding that the Einstein frame metric is asymptotically flat.
\end{enumerate}
The SL(2, R)-invariant Einstein frame metric is given by
\begin{equation}
g_{mn} \equiv e^{-\Phi/2}G_{mn}.
\label{sl2r invariant Einstein frame metric}
\end{equation}
From string perturbation theory, the string frame metric $G_{mn}$ is taken to be asymptotically flat: $G_{mn} \to \eta_{mn}$, hence the SL(2, R)-invariant Einstein frame metric is not: $g_{mn} \to g_s^{1/2}\eta_{mn}$.
There is an implicit difference in length scale between the string and Einstein frame $l_E = g_s^{1/4}l_s$, which one must be careful when comparing tensions computed in the Einstein and string frame.
The action is given as
\begin{equation}
\begin{aligned}
S_{IIB}
&=\frac{1}{2\kappa^2} \int d^{10}x\sqrt{-g}\left[\left(R-\frac{1}{2}\frac{\partial\bar\tau\partial\tau}{\tau_2^2}\right)
-\frac{1}{2}\left(
e^{-\Phi}|H_3|^2+e^\Phi|\tilde{F}_3|^2
\right)
-\frac{1}{4}|\tilde F_5|^2
\right]\\
&-\frac{1}{4\kappa^2}\int C_4\wedge H_3\wedge F_3.
\end{aligned}
\end{equation}
Had one chosen to taken route 2) and intend to work with an asymptotically-flat Einstein frame metric $\tilde{g}_{mn} \equiv g_s^{1/2}e^{-\Phi/2}G_{mn}$, there is no longer a difference in length scale between the two frames, but the metric is not SL(2, R)-invariant:
\begin{equation}
\tilde{g}_{mn} \to |c\langle\tau\rangle + d| \tilde{g}_{mn}.
\end{equation}
Hence to write down an SL(2, R)-invariant action, appropriate factors of $g_s$ must be included. For example, while $R(\tilde{g}_{mn})$ and $\sqrt{-\tilde{g}}\,R(\tilde{g}_{mn})$ are not SL(2, R)-invariant, 
$g_s^{1/2}R(\tilde{g}_{mn})$ and $\frac{1}{g_s^2}\sqrt{-\tilde{g}}\,R(\tilde{g}_{mn})$ are invariant.
Using the asymptotically-flat Einstein frame metric $\tilde{g}_{mn} = g_s^{1/2}e^{-\Phi/2}G_{mn}$, the type IIB action is
\begin{equation}
\begin{aligned}
S_{IIB}
&=\frac{1}{2\kappa^2} \int d^{10}x\sqrt{-\tilde{g}}\left[
\frac{1}{g_s^2}\left(R-\frac{1}{2}\frac{\partial\bar\tau\partial\tau}{\tau_2^2}\right)
-\frac{1}{2 g_s}\left(
e^{-\Phi}|H_3|^2+e^\Phi|\tilde{F}_3|^2
\right)
-\frac{1}{4}|\tilde F_5|^2
\right]\\
&-\frac{1}{4\kappa^2}\int C_4\wedge H_3\wedge F_3,
\end{aligned}
\label{asymptotically flat einstein frame action}
\end{equation}
where the factors of $g_s$ are fixed by demanding that the action is SL(2, R)-invariant.

Although the scalars $\Phi, C$, and the form fields $C_2, B_2$ have different NSNS, RR origins in 10d, they transform into each other under SL(2, R). 
This SL(2, R)-invariance will hold at the two-derivative level, as well as higher-derivative corrections with trivial dependence on $\tau$.
But as soon as any corrections to type IIB supergravity with non-trivial dependence on $\tau$ enter, the SL(2, R) symmetry will be broken.
Then as one includes the corrections that account for the string and brane-effects, the symmetry will also become SL(2, Z).
This reflects the quantization of NSNS and RR charges in type IIB string theory.
}

{\color{red}
As we have discussed earlier, the NSNS vs RR distinction is really a 10d one.
Both the NSNS and RR fields in type IIA combine to form SO(9) multiplets.
Analogously, the NSNS and RR fields in type IIB combine to form $SL(2, R)\times SO(8)$ multiplets, and it is natural to ask whether this could be a hint of a higher-dimensional origin?
This is the third puzzle of the type IIB theory, and it is unlikely to be independent of the previous two: where does the SL(2, Z)\footnote{
SL(2, R) at the two-derivative supergravity level.
} symmetry come from\footnote{
The type IIB theory is understood as the decompactifying limit of M-theory on a torus, which could explain SL(2, Z). But SL(2, Z) is a true symmetry in 10d already.}?
We approach this review motivated by the 3 guiding questions introduced above:
\begin{enumerate}
    \item Does the type IIB theory have a higher-dimensional origin?
    \item What is the underlying mechanism of $F_5 = *_{10}F_5$?
    \item What is the role and origin of SL(2, Z)?
\end{enumerate}
}